\chapter*{Prefácio}
\addcontentsline{toc}{chapter}{Prefácio}
\markboth{Prefácio}{Prefácio}

O tema desta tese é o modo como se faz pesquisa em inteligência artificial em uma universidade brasileira, no período em que a inteligência artificial generativa se tornou pública e deixou de ser assunto restrito a especialistas. O seu objeto é o Centro de Inteligência Artificial da Universidade de São Paulo (C4AI/USP), criado em 2020 por meio de um edital público e um convênio entre essa universidade, a fundação estadual de amparo à pesquisa de São Paulo (FAPESP) e a corporação International Business Machines (IBM). Esse convênio não foi renovado no início de 2026, e com ele terminou a parceria da corporação com o Centro, que continua existindo, apoiado pela fundação e pela universidade. Sigo cientistas e engenheiros do Centro entre março de 2022 e dezembro de 2025: entre eles, Fábio e Cláudio no capítulo~\ref{capitulo3}, e Marcelo e o Spira, um sistema de detecção de insuficiência respiratória pela voz, no capítulo~\ref{capitulo4}. O meu objetivo é descrever modos situados de fazer pesquisa em inteligência artificial. A metodologia é seguir os atores universidade afora, ou os actantes, termo que prefiro e cuja distinção esclareço adiante. Esta tese se guia por uma pergunta, a de onde está o laboratório de inteligência artificial e onde estão os seus cientistas, que desdobro em como cientistas e engenheiros fazem inteligência artificial, e por que isso importa, para quê e para quem. E a conclusão a que chego é que o laboratório de inteligência artificial se distribui, entre computadores pessoais, reuniões virtuais e infraestruturas remotas compartilhadas, enquanto a infraestrutura que sustenta esse trabalho se concentra nas mãos de poucos atores no mundo, o que torna mais difícil construir alternativas de inteligência artificial fora deles; e que a inteligência artificial generativa é feita e refeita a cada vez, distinta conforme quem a interpela e o que faz com ela. Faço esta tese com uma inteligência artificial de uma dessas poucas empresas. Escolho \textit{ficar com o problema}, permanecer nessa contradição e tensioná-la ao longo da pesquisa, sem a pretensão de resolvê-la nem de escapar dela recusando-a ou aceitando-a acriticamente, e é dela que parte a provocação que deixo em aberto: o que significa fazer ciência, e as ciências sociais em particular, com esses tipos de inteligências artificiais, e por quê. As respostas que ofereço são, por isso, provisórias e incompletas.

Intitulo esta tese de \textit{tecnografia}. Mas o que é uma tecnografia? A etnografia sempre se fez de técnicas variadas. Com o lápis e a caneta se fazem a escrita, os desenhos, os mapas, os rabiscos e os traços; com a câmera, as fotografias; com o gravador, o registro da fala; com a máquina de datilografia se imprimiam letras, números e símbolos no papel; com o computador e os \textit{softwares} se fazem textos, imagens e vídeos digitais; e cada caderno de campo, cada gráfico, cada anotação é o próprio \textit{campo} sendo inscrito. Em um sentido mais estrito, o da antropologia da técnica, chamo de técnica as operações que transformam a matéria e de tecnologia o sistema que elas formam. Em um sentido mais amplo, o dos estudos de ciência e tecnologia, reúno as duas, com a ciência que as atravessa, sob o nome de \textit{tecnociência}, e é esse o termo que uso ao longo da tese quando não separo o fazer científico do técnico. \textit{Tecnografia} une \textit{tecno} e \textit{grafia}. A raiz de \textit{grafia}, o grego \textit{gráphein}, é ao mesmo tempo escrever e traçar, desenhar, riscar, e por isso traduzir \textit{grafia} apenas por escrita a estreita demais. A grafia é o texto, a imagem, o vídeo, e também o código e a linguagem de programação, o diagrama, o gráfico, o espectrograma, o site da tese, e é nesse sentido que a tecnografia é a grafia da técnica, a técnica tomada como aquilo que descrevo, e a grafia com técnicas, a técnica tomada como o modo pelo qual a tese se faz. O meio com que pesquiso é parte do que a pesquisa diz, porque cada técnica medeia e transforma o que passa por ela, e o que assim se transforma entra no argumento em vez de apenas transmiti-lo. É por isso que digo que a tecnografia reúne em um só gesto a teoria, o método e a análise: a grafia com que faço a tese é também aquilo que a tese sustenta. As tecnografias são, portanto, os modos situados pelos quais esta tese é feita, e por fazer entendo a observação, a interpretação, a descrição, a apresentação e a análise que compõem essa feitura.

A inteligência artificial é aquilo que o Centro pesquisa, e se tornou o meu objeto de estudo à medida que eu acompanhava como os cientistas e engenheiros a fazem, e passou também a ser uma das tecnografias com que faço a tese. A inteligência artificial reúne muitas máquinas, conhecimentos de várias disciplinas, dados acumulados, técnicas de treinamento e aplicações que se multiplicam, e a tese a trata como um modo de conhecimento, feito de práticas distribuídas e postas em relação a partir deles. Esse modo carrega as suas questões, os seus efeitos e os seus impactos, da concentração em poucas empresas ao custo material, ao trabalho precarizado, aos vieses e à opacidade. Considero essas questões onde a pesquisa as encontra, no que descrevo em campo e no que discuto com a teoria, ao lado da literatura crítica que as sustenta. Essas questões estavam sendo levantadas ao mesmo tempo pelos especialistas que fazem inteligência artificial e pelas pessoas, instituições de ensino e profissões que começavam a ser por ela afetadas, e foi uma curiosidade antropológica sobre como a inteligência artificial generativa funciona na prática que me levou a experimentá-la, e foi assim que ela entrou na feitura desta tese. Foi por esse caminho que passei a entendê-la também como um modo de conhecimento que acolho sob o conceito de tecnografia.

Esta tese é, ela mesma, uma tecnografia. Escrevi-a no Overleaf, um editor de textos em \textit{LaTeX} no navegador de internet, ligado ao GitHub, um serviço que guarda e versiona os arquivos em um repositório público, e é dele que se publica o site da tese. O texto que você tem em mãos é a parte mais visível de uma cadeia que continua nos programas, nos dados e nas inscrições que ficaram guardados e abertos, e que você pode conferir pelo site e pelo repositório.

A inteligência artificial é parte dessa feitura, de vários modos. Com o Claude Code, ferramenta de codificação baseada em agente de inteligência artificial da Anthropic, trabalhei para entender e operar os softwares com os quais fiz a tese, cada um com a sua linguagem.\footnote{O Claude e o Claude Code são duas inteligências artificiais da empresa Anthropic que dão acesso ao mesmo modelo de linguagem por interfaces distintas: o Claude é a interface de conversa, em janela de chat no navegador de internet ou em aplicativo, e o Claude Code é a interface de linha de comando, executada no terminal, com acesso ao sistema de arquivos para ler, escrever e executar \textit{scripts}.} Alguns desses softwares são inteligências artificiais, como o próprio Claude; outros trazem inteligência artificial integrada, como o Overleaf, o GitHub e o Mermaid, com que gerei os diagramas em arquivos de texto; e alguns, como o GitHub e o site da tese, eu sincronizei com o Claude Code. Com o Claude Code também busquei os conjuntos de dados bibliográficos das plataformas SciELO, CAPES e OpenAlex, a partir dos quais produzi os gráficos das análises bibliométricas, examinei a rede textual da tese e fiz a lexicometria das figurações nas obras de Latour para representá-la graficamente. As análises correram em Python, reproduzido no Anaconda, um ambiente que reúne essa linguagem e as suas bibliotecas, e o Claude e o Claude Code instalados no meu computador foram o que me permitiu baixar e analisar conjuntos de dados muito grandes, como os dessas três plataformas. Com o Claude, nas versões Sonnet e Opus, na sua janela de conversa, trabalhei na revisão de ortografia e gramática, na organização da estrutura da tese em capítulos, seções e subseções, nos resumos, na tradução, na organização das referências bibliográficas em um arquivo \textit{.bib}, e em uma revisão recursiva das minhas questões teóricas e metodológicas, que me serviu para identificar lacunas, contradições, repetições e incoerências do argumento, e também os seus pontos mais bem fundamentados. Fora do Claude, usei o NotebookLM, do Google, para organizar conjuntos temáticos de obras e buscar dentro deles onde e como uma ideia se desenvolve, em que datas e páginas aparece, com que frequência e com quais outros autores dialoga. Tudo isso se deu dentro dos limites do que a inteligência artificial pode fazer, e a partir de condições que preparei: conjuntos de dados que curei com cuidado, um entendimento suficiente dos processos estatísticos e das suas automações, e uma bibliografia que eu já conhecia bem. Foi a inteligência artificial que tornou possível, nesses limites, aprender a interagir com essas plataformas e organizar, descrever, analisar e representar visualmente os dados, as informações e o conhecimento que a pesquisa produziu.

Os diagramas, os gráficos, os códigos e as análises estatísticas que compõem o texto são o argumento se fazendo. Isso é mais saliente no capítulo~\ref{capitulo4}, sobre o Spira, um sistema que detecta insuficiência respiratória pela análise da voz. Ali trabalho com o mesmo material dos artigos científicos do Spira, as vozes das pessoas que contribuíram para desenvolver o sistema, que busquei no conjunto de dados público do Spira no GitHub e converti em representações visuais, os espectrogramas, junto ao Claude Code. A equipe extrai inscrições tecnocientíficas desses espectrogramas publicados nos artigos, e eu tomo essas mesmas vozes e as reinterpreto, produzindo com elas inscrições tecnográficas. A distinção entre as duas eu construo ali, e o que a marca é o objetivo: nos artigos a voz se torna dado e se apaga como som, e o que faço é ouvir uma dessas vozes e dá-la a ouvir, inserindo o seu áudio na tese e no site para que você também a escute. É esse o sentido tecnográfico que estendo à voz: o espectrograma que produzo junto ao Claude Code passa, ao entrar na tese, de objeto da análise a inscrição tecnográfica com que a faço.

Eis por que digo que \enquote{faço com} --- \textit{fazer-com}.

Fazer é sempre fazer com algo, alguém ou alguma coisa. O campo, os interlocutores, a orientação e as técnicas com que trabalho fizeram esta tese comigo. No \textit{fazer-com}, quem faz e aquilo com que se faz não estão prontos antes da relação, e se constituem um ao outro no próprio fazer. É assim que trabalho com a inteligência artificial e com tudo o mais que pus em relação nesta tese. Teria feito esta tese sem a inteligência artificial, mas a tese que fiz com ela é outra, porque o \textit{fazer-com} deixa a sua marca no que se faz, e é essa marca que dá maior ênfase à tecnografia. Chamo de tecnografia o que faço, e nela incorporo os modos de conhecimento e as expressões das tecnociências, as inscrições que o campo produz, os gráficos, os espectrogramas, as tabelas e os códigos, postos em relação com a escrita, com a teoria e com o campo. A tecnografia é, por isso, um conhecimento relacional, que resulta de pôr essas coisas umas em relação às outras, e um conhecimento situado, que faço de um lugar preciso, o das ciências sociais, como pesquisadora sobre a inteligência artificial e com a inteligência artificial. O que se faz aí não é nem tão somente meu nem tão somente das técnicas, é da relação entre nós --- um \textit{híbrido}, se se quiser pensar dessa maneira. A ação se distribui pela rede, e com ela a responsabilidade pelo que a inteligência artificial faz se dilui entre todos os que fazem essa rede, as empresas que a produzem, quem as desenvolve, quem recorre a elas, eu. Repartir essa responsabilidade com precisão é impossível, e na prática ela recai sobre quem faz a pesquisa com a inteligência artificial, o que não é o ideal nem o mais justo. 


No meu caso, sabendo dos problemas e dos limites da inteligência artificial, escolhi fazer a minha tese com ela, atenta às recomendações do regimento e aos limites éticos, técnicos e científicos da Universidade,\footnote{O uso de ferramentas de inteligência artificial generativa nesta pesquisa observa a Deliberação CONSU-A-005/2026 da Universidade Estadual de Campinas, e está detalhado na Declaração de Uso de Inteligência Artificial Generativa, no Anexo~1.} e é essa escolha que faz a responsabilidade distribuída pesar mais sobre mim, que respondo pela tese e por tudo o que nela há. Esse híbrido é um novo modo de produção de conhecimento, cujos efeitos ainda estão sendo investigados, e friso isso porque é o que dá a medida do desafio a que me proponho, diante de um contexto que vai do deslumbramento ao medo das inteligências artificiais. Para mostrar como esse \textit{fazer-com} se deu, disponibilizo o processo tanto quanto é possível: os repositórios versionados, os códigos e os passos de cada análise ficam públicos, abertos a quem queira refazê-los ou discordar deles.

Quem são os atores fundamentais que você vai encontrar nesta tese? A universidade, a fundação e a corporação sustentaram o Centro, que reuniu, em seu período de maior atividade, cerca de 250 pesquisadores espalhados por vários campi e instituições. Sigo Fábio Cozman, diretor do Centro, e Cláudio Pinhanez, vice-diretor e, na época, pesquisador da corporação, pela rede do Centro no capítulo~\ref{capitulo3}, e sigo Marcelo Finger, coordenador do grupo de processamento de linguagem natural para o português e coordenador do projeto do Spira, pela rede desse sistema no capítulo~\ref{capitulo4}.

Digo atores por conveniência, mas o termo que adoto é \emph{actante}, que prefiro a ator para não supor que agir dependa de intenção ou de consciência humana. Chamo de actante qualquer coisa, gente ou não, que muda outra no curso de uma relação, e um contrato, um computador, uma voz gravada na enfermaria, o vírus da Covid-19 e uma linguagem de programação são actantes na medida em que desviam o rumo das relações em que participam. No campo, as redes que sigo misturam os dois antes de qualquer divisão analítica. Na escrita, porém, distingo humano e não humano como um recurso de interpretação e descrição, que me serve para recusar a divisão prévia entre o sujeito que age e o objeto que só se deixa usar. Esse recurso deriva das tecnociências ocidentais e da própria literatura dos estudos de ciência e tecnologia, que hoje se desloca do par humano e não humano para formulações como a do além do humano.

Esses actantes foram fundamentais porque foi a partir deles que consegui identificar o maior número de associações: as que ligam a universidade, a fundação e a corporação, com os seus convênios e financiamentos; as que ligam os pesquisadores às inteligências artificiais e aos dados; e as que ligam a voz de um paciente ao dado, ao artigo e ao algoritmo, e a parceria de hoje à longa história da tabulação de dados. O subtítulo desta tese fala em seguir cientistas e engenheiros, e segui-los me levou a seguir também os computadores, as inteligências artificiais, os dados, os contratos e as vozes que eles reúnem, porque é disso que as redes são feitas. Considerei essas associações como os nós mais densos da rede. E as redes que sigo são sociotécnicas e tecnocientíficas, associações entre actantes heterogêneos, humanos e não humanos, em que o social e o técnico, a ciência e a técnica se fazem juntos.

Além da observação e da imersão nas atividades cotidianas do Centro, realizei entrevistas com os seus pesquisadores. A realização da pesquisa foi autorizada por escrito pelo Centro, e cada entrevista, feita por videoconferência, foi autorizada pelos interlocutores por escrito e oralmente. Por ser esta uma etnografia que segue os atores, as entrevistas aconteceram conforme os interlocutores apareciam na rede que eu ia acompanhando. Cada uma seguiu um roteiro semiestruturado, com perguntas sobre a trajetória acadêmica e o trabalho desenvolvido no Centro, que mantive flexível para deixar a resposta o mais aberta e livre possível, de modo que cada interlocutor desse à sua fala o rumo que quisesse. Entrevistei os coordenadores dos cinco grupos de pesquisa do Centro, e trago para o centro três deles, Fábio, Cláudio e Marcelo, porque eram os pontos a partir dos quais a rede se deixava seguir com mais facilidade e porque as suas falas abriam  caminhos para os dois principais argumentos desta tese, as relações institucionais no capítulo terceiro e o do sistema sociotécnico no quarto. Os demais pesquisadores e estudantes aparecem por outros caminhos, as conversas de campo, os relatórios, as publicações. Gravei e transcrevi as entrevistas para voltar a elas na análise, e as guardo nos meus registros de pesquisa. Escolhi trabalhar as falas de modo contextual e situado, no corpo do texto, em vez de anexar as entrevistas na íntegra, por questões éticas. 

Quais materiais fazem parte desta pesquisa? Entrei em campo em março de 2022, e reconstruí o período anterior, da criação do Centro à minha entrada, a partir de relatórios e entrevistas, e é como reconstrução que ele entra aqui. Boa parte da tese trabalha com o passado, e sigo cadeias que recuam no tempo por documentos, registros, contratos e acordos, editais públicos, sites institucionais, videoaulas, obras biográficas e históricas e artigos científicos. A esses materiais que reúno se soma o que eu mesma inscrevo, as descrições do meu caderno de campo, onde registrei a observação entre março de 2022 e dezembro de 2025. Seguir a rede, inclusive a que as próprias entrevistas iam indicando, com as suas associações particulares, me levou a esses materiais e a esses modos de análise. A pesquisa se fez, assim, multissituada, acompanhando os atores por muitos lugares ao mesmo tempo, a universidade, os laboratórios, os eventos públicos, os repositórios; e se fez também de arquivo, ao analisar conjuntos de dados de inteligência artificial, a bibliometria do campo, os registros que o Centro deixou e a trajetória da corporação, que reconstruo no capítulo terceiro a partir dos seus documentos institucionais e da história da tecnologia. É dessa junção que vem uma dimensão particular da tese, a de integrar a observação no presente --- em tempo real --- e a análise do passado --- em retrospecto. Por isso o que infiro da observação, das entrevistas e da análise desses materiais é mais um trabalho de fazer relações do que de diagnóstico e explicação.

Como ler os capítulos desta tese? Eles têm caráter distinto. O primeiro é metateórico, e se volta sobre as próprias grafias da tese. O segundo é uma análise bibliométrica do campo, lexicométrica de obras selecionadas e também uma reflexão sobre o papel das figurações nos estudos das tecnociências. O terceiro segue a rede que Fábio e Cláudio construíram, e o quarto, a rede que Marcelo construiu. Cada um foi escrito em um momento distinto da pesquisa, que coincidiu com a chegada pública da inteligência artificial generativa, e isso abriu de uma só vez uma variedade de questões com as quais tive de lidar. Daí vem a heterogeneidade dos capítulos, cada um respondeu às questões de uma fase da pesquisa.

Começo pelo metateórico e pela análise bibliométrica, e deixo para depois as redes de Fábio, Cláudio e Marcelo. Os dois primeiros poderiam parecer os menos etnográficos, e a ordem, um percurso que vai da teoria ao campo. Não é o caso. A etnografia atravessa toda a tese, e está também nos capítulos mais teóricos, cuja teoria surgiu das idas a campo, de modo que o que parece reflexão abstrata é o registro escrito de um percurso concreto. Começo por eles porque preparam o vocabulário e as decisões com que sigo o Centro e o Spira, e é por já terem feito esse trabalho que o terceiro e o quarto seguem essas redes com mais leveza. Você pode, ainda assim, entrar pelo terceiro e pelo quarto e voltar aos primeiros quando quiser reencontrar o que os sustenta.

Cada capítulo abre com dois diagramas, um mapa e uma linha do tempo. Chamo o primeiro de mapa porque é um tipo de guia: dá a ver por onde o capítulo começa e aonde chega, da pergunta que o abre às suas seções, com as caixas e as setas de que é feito. O mapa é outra tecnografia da tese e apresenta as ideias e o argumento do capítulo em forma visual e resumida. Pelo tamanho, cada mapa traz na legenda um \emph{link} para uma versão interativa, e você pode seguir a leitura sem abri-lo, ou abri-lo quando quiser ver o capítulo de outro modo.

\begin{center}
  * * *
\end{center}
A tese também pode ser lida como rede textual na internet.\begin{center}
  {\ttfamily\small\color{latexneon}\textbackslash url\{%
   \color{black}\href{https://julianehelanski.github.io/tecno-etnografia-tese-site/}%
   {julianehelanski.github.io/tecno-etnografia-tese-site}\color{latexneon}\}}\par
  \vspace{4pt}
  {\small\texttt{Para uma melhor experiência, recomenda-se o acesso pelo computador}.}
\end{center}

